\chapter{Transesophageal Echocardiography}

\section*{Overview}
Transesophageal echocardiography (TEE) is an ultrasound examination of the heart performed with a probe passed into the esophagus and stomach, allowing imaging of the heart from directly behind it. It provides high-resolution images of the cardiac structures, valves, and great vessels.

\section*{Alternative Names}
TEE; esophageal echocardiogram; transesophageal echo

\section*{Principle}
A transducer at the tip of a flexible endoscope emits and receives ultrasound waves from within the esophagus, which lies immediately posterior to the left atrium and aorta. The short distance between the probe and the heart avoids interference from the chest wall and lungs, producing superior image quality compared with transthoracic imaging.

\section*{History}
Transesophageal echocardiography was developed in the 1970s and entered widespread clinical use in the 1980s to 1990s, becoming a standard tool for evaluating structures that are difficult to visualize from the chest wall.

\section*{Why This Test or Procedure Is Ordered}
Ordered to evaluate infective endocarditis and its complications, atrial thrombi before cardioversion or ablation, prosthetic valve function, the source of cardioembolic stroke, aortic dissection, congenital heart defects, and to guide cardiac surgery and interventional procedures.

\section*{Is This a Primary or Secondary Test or Procedure}
Secondary to transthoracic echocardiography for most indications; it is the primary modality for infective endocarditis assessment, prosthetic valve evaluation, and intraoperative cardiac monitoring.

\section*{How This Test or Procedure Is Performed}
The patient fasts and is placed in the left lateral position after topical anesthesia and conscious sedation. A lubricated transesophageal probe is passed into the esophagus under guidance, and multiple two- and three-dimensional images, Doppler measurements, and, when indicated, three-dimensional reconstructions are acquired.

\section*{What Preparation Is Needed}
Fasting for at least four to six hours, informed consent, review of swallowing disorders and anticoagulation, and removal of dentures are needed. The patient should arrange for someone to accompany them home after sedation.

\section*{Contraindications and Precautions}
Esophageal stricture, esophageal or gastric varices, recent esophageal or gastric surgery, esophageal perforation, and active upper gastrointestinal bleeding are contraindications. Precautions include monitoring oxygen saturation and blood pressure during the procedure and avoiding the study in patients with suspected esophageal pathology.

\section*{Risks}
Risks include transient throat discomfort, mucosal injury, and rarely esophageal perforation, bleeding, arrhythmia, or respiratory compromise from sedation.

\section*{What Happens After the Test or Procedure}
The patient is observed until sedation wears off, and the throat may be sore for a short period. Food and drink are resumed after sensation returns and the patient is confirmed to be able to swallow safely.

\section*{Understanding Your Results}
Findings such as valvular regurgitation, vegetations, thrombus, or dissection are interpreted by the cardiologist and integrated with clinical and laboratory data. Abnormalities guide decisions on surgery, anticoagulation, or further imaging.

\section*{Normal Results}
Normal findings show intact cardiac structures, normally functioning valves without vegetations or regurgitation, no thrombus, and normal chamber sizes and ventricular function.

\section*{Follow-up Tests and Procedures}
Abnormal findings may lead to cardiac surgery, repeat imaging, or blood cultures in suspected endocarditis. Follow-up TEE is used to monitor prosthetic valves, assess therapy response, or reevaluate thrombus resolution before cardioversion.

\section*{References}
\begin{enumerate}
  \item MedlinePlus. Transesophageal echocardiography (TEE). U.S. National Library of Medicine; 2025.
  \item Current Medical Diagnosis and Treatment. McGraw-Hill; 2025.
\end{enumerate}

\clearpage
